\section{Conclusion and Future Work}
\label{sec:conclusion}

\subsection{Summary of contributions}
TriageGuard turns a Windows or Linux memory image into numbered, cited findings, a sourced risk score with an
uncertainty range, and a readable report that an LLM writes but cannot be steered by the evidence. Its
cryptographic stage goes past naming ciphers: it states the evidence basis of each identification, grades
every RSA key against NIST, and recovers live AES, ChaCha20 and RSA private keys, including four RSA-2048 keys
from real WannaCry memory. On 14 images the scores separate a clean baseline (28) from compromised hosts (60 to
97); all 12 ground-truth checks pass; and all 14 LLM reports cite 90 to 100\% of their claims with the risk stated
exactly. Across 126 defended red-team reports, including adaptive attacks, no attack succeeded and none would
have misled a reader without a warning. All seven objectives are met.

Three lessons generalise. Real data, not synthetic samples, drove most of the useful changes: the first
Windows~11 run produced 689 findings and a risk of 100, and every benign pattern we later downgraded came from
reading real output. Identification by constants overstates what is in use, so recovered keys are the stronger
evidence. And input filtering alone does not stop an adaptive attacker; checking the output against
deterministic results does.

\subsection{Limitations}
\begin{itemize}
  \item Linux images need the exact kernel's symbol table; without it Volatility cannot parse the image.
  \item Keys are recovered only while live (0/15 after release in the benchmark), except where the platform
        leaves them behind, as legacy CryptoAPI does.
  \item Signature scanning finds known constants; a custom or obfuscated cipher would not match.
  \item The RSA prime search skips processes over 192\,MB; the pcap reader stops at 50{,}000 packets.
  \item Red-team cells hold 6 trials, one model and synthetic payloads, so 0/126 bounds the unflagged rate at
        about 3\% rather than proving it zero.
  \item Campaign linking produces false links between images of the same Windows build.
  \item Severities and region grading are heuristics tuned on public images.
\end{itemize}

\subsection{Future work}
\begin{itemize}
  \item \textbf{Linking baseline.} Learn the known-normal process and endpoint set from clean images of each OS
        build, so stock apps and vendor endpoints stop creating links.
  \item \textbf{More recovery targets.} TLS session secrets~\cite{baier2024tls}, Salsa20 ransomware
        contexts~\cite{babiano2023evaluation} and elliptic-curve private scalars.
  \item \textbf{Stronger red-team.} More trials per cell, several models (including a local one), multi-field
        attacks and attacks aimed at the cross-check itself.
  \item \textbf{Post-quantum readiness.} ML-KEM (FIPS~203) and ML-DSA keys will start to appear
        in memory; identifying their constants and flagging hosts that still rely only on RSA and ECC for key
        exchange is a natural extension of the cryptographic stage.
  \item \textbf{Scale.} Running Volatility plugins in parallel and caching symbol tables to bring the 22-minute
        run on a 4.3\,GB image down for larger servers.
\end{itemize}

\section*{Code and data availability}
Source code, all exported results and the formulas behind every score are at \repo; a read-only dashboard over
all 14 images is at \demo. Evaluation used public training images and synthetic keys only.
